% Ссылка в тексте: (приложение~Д)
\chapter{СПЕЦИФИКАЦИЯ ГРАФА АГЕНТОВ LANGGRAPH}

В~приложении представлена текстовая спецификация фактического
шестиузлового графа агентов ИИ-методиста: узлы, переходы, точки участия
человека и контур отката к прежней топологии.

\section*{Д.1 --- Ноды графа}

\begin{table}[h!]
  \caption{Ноды графа агентов ИИ-методиста}
  \label{tab:langgraph-nodes}
  \centering
  \begin{tabular}{|p{3.5cm}|p{2.5cm}|p{5.5cm}|}
  \hline
  \textbf{Нода} & \textbf{Тип} & \textbf{Описание} \\
  \hline
  source & агент & Поиск источников, EvidenceTable, статус покрытия \\
  \hline
  structure & агент & CourseBlueprint по методологии ADDIE и доказательному контексту \\
  \hline
  HITL-structure & прерывание & Подтверждение или коррекция структуры курса \\
  \hline
  content & агент & Черновики уроков со встроенными ссылками на источники \\
  \hline
  fact-check & агент & Проверка утверждений относительно EvidenceTable \\
  \hline
  review & агент & Методологическая проверка, покрытие Блума, связность \\
  \hline
  verifier & агент & Знак качества, подтверждённые утверждения, происхождение результата, снимок проверки дипломных метрик \\
  \hline
  HITL-quality & прерывание & Финальное решение человека перед публикацией \\
  \hline
  \end{tabular}
\end{table}

\section*{Д.2 --- Основные переходы}

\begin{table}[h!]
  \caption{Ключевые переходы в~графе LangGraph}
  \label{tab:langgraph-edges}
  \centering
  \begin{tabular}{|p{3.2cm}|p{3.2cm}|p{5cm}|}
  \hline
  \textbf{Откуда} & \textbf{Куда} & \textbf{Условие перехода} \\
  \hline
  source & structure & EvidenceTable построена или явно помечена как неполная \\
  \hline
  structure & HITL-structure & структура курса доступна для проверки \\
  \hline
  HITL-structure & content & структура утверждена пользователем \\
  \hline
  HITL-structure & structure & структура отклонена и требует пересборки \\
  \hline
  content & fact-check & уроки подготовлены и требуют верификации \\
  \hline
  fact-check & review & отчёты проверки фактов сформированы \\
  \hline
  review & content & требуется доработка, revision count меньше лимита \\
  \hline
  review & verifier & курс прошёл методологическую проверку или исчерпан лимит доработок \\
  \hline
  verifier & HITL-quality & знак качества и происхождение результата собраны \\
  \hline
  HITL-quality & END & человек принимает результат или фиксирует остановку \\
  \hline
  \end{tabular}
\end{table}

\section*{Д.3 --- Точки участия человека}

\begin{itemize}
  \item \textbf{HITL~1}: подтверждение структуры курса до генерации полного контента;
  \item \textbf{HITL~2}: финальное решение на основе знака качества,
    происхождения результата, ссылок и блокирующих замечаний;
  \item \textbf{Откат}: при отключении конвейера Feynman сервис может
    перейти к прежней четырёхузловой топологии, а происхождение результата фиксирует,
    какой конвейер фактически исполнялся.
\end{itemize}
